% page 343 du fac-simile (image p-342 du scan)
% bloc 157.5 x 250.6 mm ; marge gauche 44.5 mm
\pgc{22.860mm}{0.000mm}{
\l{}
\l{}
\l{}
\l{}
\l{}
\l{\sou{leporo}. (\sou{zool.}) Quadripedo rodera, di qua la gambi avana}
\l{\cel{3}esas plu kurta kam le dopa, e kun oreli longa. - L.\cel{1}\sou{lepus}}
\l{\cel{3}\sou{timidus}. - FIS.}
\l{}
\l{\sou{lepro}. (\sou{patol.}). Morbo di qua la fonto esas tuberkloso, qua}
\l{\cel{3}invadas, rodas e repugnante sulkizas la pelo. - DEFIS.}
\l{}
\l{\sou{lernar}. (\sou{trans.}) Aquirar la savo (di ulo) per spirito-laboro.}
\l{\cel{3}- DE}
\l{}
\l{\sou{lesivar}. (\sou{trans}.)Netigar (linjo) per lasar lu dum kelka tempo}
\l{\cel{3}imersita en aquo varma a qua onu adjuntis sodo, natroxo o}
\l{\cel{3}potaso, e qua esas apta dissolvar, ek la linjo sordida,}
\l{\cel{3}la kraso e la nepuraji qui esas an lu, adherante. - FI.}
\l{}
\l{\sou{letargio}. (\sou{patol.}) Stando maladesala quan karakterizas dormo}
\l{\cel{3}profunda e tre duroza,dum qua la fenomeni videbla di vivo}
\l{\cel{3}esas interruptita, e qua prizentas la aspekto di morto.}
\l{\cel{3}- DEFIRS.}
\l{}
\l{\sou{letro}. Skriburo quan onu sendas a persono absenta, por komuni-}
\l{\cel{3}kar a lu to quon onu ne povas dicar a lu voce. - EFIS.}
\l{}
\l{\sou{leucito}. (\sou{kemio)}. Duopla silikato di alumino e di potaso,}
\l{\cel{3}quan onu renkontras en la tereni volkan-agita. - DEFIRS.}
\l{}
\l{\sou{leukocito}. (\sou{anat.)}\cel{3}Celulo senkolora di sango e di limfo,}
\l{\cel{3}nukleoza, deformebla,e di qua la dimensioni esas plu granda}
\l{\cel{3}kam olti di hematito. - DEFIRS.}
\l{}
\l{\sou{leukomo}. Makulo blanka, opaka, qua naskas ulakaze sur la}
\l{\cel{3}korneo di la okulo. - DEFIRS.}
\l{}
\l{\sou{levar}. (\sou{trans}.)Igar movata, de infre ad supre. - EFIS.}
\l{}
\l{\sou{levero}. Stango rigida, moviva sur apogo-punto, e quan onu}
\l{\cel{3}uzas por levar, por igar objekto movar. - EFI.}
\l{}
\l{\sou{leviatano}. (\sou{biblo}). Monstro marala. - DEFIRS.}
\l{}
\l{\sou{levito}. I. (\sou{biblo}) Izraelido (ek la tribuo di Levi) konsakrita}
\l{\cel{3}a la oficio en la Templo. -II. (\sou{metaf.}) (\sou{poezio}).Sacerdoto.}
\l{\cel{3}DEFIRS.}
\l{}
\l{\sou{levrier}o.(\sou{zool.})\cel{2}Hundo kun muzelo akuta, korpo tenua, gambi}
\l{\cel{3}longa, di qua la aluro esas rapida, uzata partikulare por}
\l{\cel{3}persequar la lepori. - FIS.}
\l{}
\l{\sou{+leyo.}\cel{2}Regulo konstanta, universala, a qua omna fenomeni di}
\l{\cel{3}naturo esas obedienda. -- Rezultajo de analizar exakte}
\l{\cel{3}e precize la cirkonstanci qui esas elementi en la produkto}
\l{\cel{3}di la fenomeni, ed interligar ta cirkonstanci per la relati}
\l{\cel{3}normala di inter-sucedo e di inter-simileso. (Aug. Comte).}
}
